\section{Algorithmic and Theoretical Details}
\label{sec:appendix-algorithm-theory}

This appendix provides the technical material needed to reproduce and interpret CARAT. It first fixes the round-level notation and server inputs, then gives executable pseudocode, implementation invariants, and structural propositions for the scoring and weighting rule. These results are deliberately stated as invariants and bounds for the proposed aggregation procedure, not as a complete Byzantine-robustness guarantee against arbitrary adaptive attackers.

\subsection{Notation and Round Inputs}
\label{sec:appendix-notation}

All quantities in this appendix refer to one communication round $r$ unless stated otherwise. The server receives $n=|\mathcal{S}_r|$ client messages and converts them into vector updates $\Delta_i^r$. CARAT keeps three versions of each update:
\[
\Delta_i^r
\quad\text{(raw client update)},\qquad
\bar{\Delta}_i^r
\quad\text{(clipped update)},\qquad
\tilde{\Delta}_i^r
\quad\text{(common-radius scoring copy)}.
\]
Only $\bar{\Delta}_i^r$ is used in the final server update. The normalized copy $\tilde{\Delta}_i^r$ is used only to evaluate hidden-task utility, so the validation score is less sensitive to residual norm differences after clipping.

The server maintains a labeled reference set $\mathcal{D}^{\mathrm{ref}}$ and derives from it a working probe pool $\mathcal{P}$. Each hidden task is a small approximately class-balanced subset sampled from $\mathcal{P}$. The raw certificate matrix is $C^{\mathrm{raw}}\in\mathbb{R}^{n\times T}$, where $T$ is the number of hidden tasks and $C^{\mathrm{raw}}_{i,t}=c_{i,t}$ is the loss reduction obtained by candidate update $i$ on hidden task $t$. After robust column-wise standardization, the optimization uses $C\in\mathbb{R}^{n\times T}$. The vector $\gamma\in\mathbb{R}^n_{\ge 0}$ denotes coordinate-rank anomaly penalties. The vector $\alpha^{\mathrm{prev}}$ is the previous round's aggregation weight vector and acts as a temporal prior.

\subsection{CARAT Round Algorithm}
\label{sec:appendix-algorithm}

Algorithm~\ref{alg:carat-round} summarizes one CARAT server aggregation round. The algorithm is server-side: clients do not observe the realized hidden tasks sampled from $\mathcal{P}$.

\begin{algorithm}[t]
\caption{CARAT server aggregation at round $r$}
\label{alg:carat-round}
\begin{algorithmic}[1]
\REQUIRE Current global model $w_r$, client messages $\{u_i^r\}_{i=1}^n$, previous-round weights $\alpha^{\mathrm{prev}}$, server reference set $\mathcal{D}^{\mathrm{ref}}$, probe pool $\mathcal{P}$, CARAT hyperparameters
\ENSURE Aggregated update $\Delta_{\mathrm{agg}}$, updated temporal prior $\alpha^{\mathrm{prev}}$
\IF{$\mathcal{P}$ is uninitialized}
    \STATE Construct a class-balanced probe pool $\mathcal{P}$ from $\mathcal{D}^{\mathrm{ref}}$
\ENDIF
\IF{$r=1$}
    \STATE $\alpha^{\mathrm{prev}} \leftarrow \frac{1}{n}\mathbf{1}$
\ENDIF
\FOR{$i=1$ to $n$}
    \STATE Convert $u_i^r$ into a vector-form delta $\Delta_i^r$
\ENDFOR
\STATE Compute the clipping threshold $\tau_r$ from $\{\|\Delta_i^r\|_2\}_{i=1}^n$
\FOR{$i=1$ to $n$}
    \STATE $\bar{\Delta}_i^r \leftarrow \min\!\left(1,\frac{\tau_r}{\|\Delta_i^r\|_2+\varepsilon}\right)\Delta_i^r$
\ENDFOR
\STATE Choose the evaluation radius $q_r$ from a quantile of positive clipped norms
\FOR{$i=1$ to $n$}
    \STATE $\tilde{\Delta}_i^r \leftarrow q_r \bar{\Delta}_i^r / (\|\bar{\Delta}_i^r\|_2+\varepsilon)$
\ENDFOR
\STATE Sample $T$ hidden tasks from $\mathcal{P}$ using class-balanced sampling
\FOR{$t=1$ to $T$}
    \STATE Evaluate the baseline loss $\ell_t(w_r)$
\ENDFOR
\FOR{$i=1$ to $n$}
    \FOR{$t=1$ to $T$}
        \STATE $c_{i,t} \leftarrow \ell_t(w_r) - \ell_t(w_r + \tilde{\Delta}_i^r)$
    \ENDFOR
\ENDFOR
\STATE Robustly standardize certificate columns to obtain $C$
\STATE Compute rank penalties $\gamma$ from repeated coordinate subsamples of $\{\bar{\Delta}_i^r\}_{i=1}^n$
\STATE $\alpha^{(0)} \leftarrow \alpha^{\mathrm{prev}}$
\FOR{$k=0$ to $K-1$}
    \STATE $z \leftarrow C^\top \alpha^{(k)}$
    \STATE $\pi \leftarrow \softmin$ task weights induced by $z$ and $\beta$
    \STATE $g \leftarrow C\pi - \lambda_r \gamma - 2\lambda_p(\alpha^{(k)}-\alpha^{\mathrm{prev}})$
    \STATE $\alpha^{(k+1)} \leftarrow \Pi_{\mathcal{W}(\hat{\rho})}\!\left(\alpha^{(k)}+\eta g\right)$
\ENDFOR
\STATE $\alpha^\star \leftarrow \alpha^{(K)}$ and $\alpha^{\mathrm{prev}} \leftarrow \alpha^\star$
\STATE $\Delta_{\mathrm{agg}} \leftarrow \sum_{i=1}^n \alpha_i^\star \bar{\Delta}_i^r$
\STATE \textbf{return} $\Delta_{\mathrm{agg}}, \alpha^{\mathrm{prev}}$
\end{algorithmic}
\end{algorithm}

\subsection{Implementation Invariants}
\label{sec:appendix-implementation}

This subsection records implementation choices that affect reproducibility but are too low-level for the main method section. The purpose is to make the server-side scoring rule auditable, not to introduce additional assumptions beyond those stated in the threat model.

\paragraph{Component rationale.}
\label{sec:appendix-component-rationale}
Table~\ref{tab:component-rationale} summarizes why each CARAT component is included. The design is not intended as an arbitrary stack of heuristics: each term targets a distinct failure mode that appears when heterogeneous benign drift and malicious manipulation coexist.

\begin{table*}[!htbp]
\centering
\caption{Design rationale for the CARAT aggregation pipeline. Each component addresses a different ambiguity in non-i.i.d.\ robust aggregation.}
\label{tab:component-rationale}
\renewcommand{\arraystretch}{1.08}
\small
\resizebox{\textwidth}{!}{%
\begin{tabular}{p{0.23\textwidth}p{0.33\textwidth}p{0.36\textwidth}}
\toprule
\textbf{Component} & \textbf{Failure mode addressed} & \textbf{Role in CARAT} \\
\midrule
MAD clipping and common-radius scoring & Large-norm or scale-manipulated updates can dominate validation scores, while benign non-i.i.d.\ clients may naturally have different update norms. & Separates the scoring copy from the final clipped update, making hidden-task evaluation less sensitive to update magnitude. \\
Hidden class-balanced tasks & A single global score or trusted direction can miss class-local degradation and can confuse semantic usefulness with geometric centrality. & Scores each candidate by loss reduction on multiple hidden semantic slices, producing a task-level certificate matrix. \\
Robust task standardization & Hidden tasks can have different baseline loss scales and difficulty levels. & Makes certificate columns comparable while preserving within-task ordering. \\
Coordinate-rank penalty & A client may look useful on part of the reference signal while remaining structurally extreme in update space. & Penalizes consistently outlying coordinate ranks without rewarding unusual centrality. \\
Capped-simplex weights & Even a high-scoring client should not receive unbounded influence in a Byzantine setting. & Limits per-client weight according to an attacker-fraction prior and bounds direct malicious mass. \\
Temporal regularization & Round-to-round weight oscillations can overreact to noisy probe evaluations. & Encourages stable trust allocation unless hidden-task evidence supports a change. \\
\bottomrule
\end{tabular}
}
\end{table*}

\paragraph{Persistent state and initialization.}
CARAT maintains two state variables across rounds: the probe pool $\mathcal{P}$ derived from $\mathcal{D}^{\mathrm{ref}}$, and the previous-round weight vector $\alpha^{\mathrm{prev}}$. At the first round, $\alpha^{\mathrm{prev}}$ is initialized to the uniform vector. This makes the first aggregation unbiased across clients while allowing later rounds to use temporal consistency once weight trajectories have been observed.

\paragraph{Reference-pool construction and hidden-task sampling.}
Let $\mathcal{P}_c \subseteq \mathcal{P}$ denote the subset of probe examples with class label $c$. A hidden task is constructed by sampling a small number of examples from each available class pool and concatenating them into one balanced micro-task. If the desired number of samples per class is not available without replacement, CARAT falls back to replacement and, if needed, to a random subset of the probe pool rather than aborting the round. This fail-soft design preserves the intended semantic diversity of the task sampler while ensuring that the server always obtains a well-defined certificate matrix.

\paragraph{Robust score calibration.}
The raw certificate matrix is transformed column-wise before optimization:
\[
C_{i,t}
=
\frac{C^{\mathrm{raw}}_{i,t}-\operatorname{median}_j(C^{\mathrm{raw}}_{j,t})}
{1.4826 \cdot \operatorname{MAD}_j(C^{\mathrm{raw}}_{j,t})+\varepsilon}.
\]
This removes task-specific scale effects while preserving within-task ordering by Proposition~\ref{prop:task-standardization}. Rank penalties are calibrated similarly. For each coordinate subsample $S_s$, let
\[
\rho_i^{(s)}=
\frac{1}{|S_s|}\sum_{j \in S_s}
\left|
\frac{\operatorname{rank}_{i,j}}{n-1} - \frac{1}{2}
\right|.
\]
CARAT then computes a standardized rank score
\[
z_i^{(s)}=
\frac{\rho_i^{(s)}-\operatorname{median}_\ell(\rho_\ell^{(s)})}
{1.4826 \cdot \operatorname{MAD}_\ell(\rho_\ell^{(s)})+\varepsilon},
\]
and aggregates subsamples through
\[
\gamma_i=\max\left(0,\operatorname{median}_s z_i^{(s)}\right).
\]
Thus a client is penalized only when it is consistently extreme across repeated coordinate views, not because of a single noisy subsample.

\paragraph{Projected optimization.}
CARAT initializes the projected optimizer with $\alpha^{(0)}=\alpha^{\mathrm{prev}}$. If
\[
z(\alpha)=C^\top \alpha,
\qquad
\pi_t(\alpha)=\frac{\exp(-\beta z_t(\alpha))}{\sum_{\ell=1}^T \exp(-\beta z_\ell(\alpha))},
\]
then the ascent direction used in projected updates is
\[
g(\alpha)=C\pi(\alpha)-\lambda_r \gamma-2\lambda_p(\alpha-\alpha^{\mathrm{prev}}).
\]
With optimizer step size $\eta$, one iteration takes the form
\[
\alpha^{(k+1)}=
\Pi_{\mathcal{W}(\hat{\rho})}
\left(\alpha^{(k)}+\eta g(\alpha^{(k)})\right).
\]
The projection maintains feasibility throughout optimization, so every intermediate iterate satisfies nonnegativity, unit-sum, and per-client cap constraints.

\paragraph{Diagnostics.}
The implementation records learned weights, task scores, softmin task weights, rank penalties, clipping thresholds, evaluation radii, and hidden-task identifiers. These diagnostics distinguish three qualitatively different behaviors: benign but diverse clients that retain moderate weight, harmful clients that lose certificate support across hidden tasks, and structurally extreme clients whose influence is reduced mainly through the rank penalty.

\subsection{Proofs and Structural Properties}
\label{sec:appendix-theory}

The following propositions are grouped by the role they play in CARAT. Propositions~\ref{prop:capped-simplex}--\ref{prop:malicious-mass-bound} characterize feasible weighted aggregation and bounded influence after clipping. Propositions~\ref{prop:directional-invariance}--\ref{prop:task-standardization} characterize the hidden-task evaluation geometry. Propositions~\ref{prop:softmin}--\ref{prop:pairwise-transfer} characterize the optimization objective, and Proposition~\ref{prop:complexity} records the per-round cost.

\begin{proposition}[Feasibility and support of the capped simplex]
\label{prop:capped-simplex}
Let
\[
\mathcal{W}(\hat{\rho})=
\left\{
\alpha \in \R^n:
\alpha_i \ge 0,\;
\sum_{i=1}^n \alpha_i = 1,\;
\alpha_i \le \frac{1}{(1-\hat{\rho})n}
\right\},
\]
with $\hat{\rho}<1$. Then $\mathcal{W}(\hat{\rho})$ is non-empty, and every $\alpha \in \mathcal{W}(\hat{\rho})$ satisfies
\[
|\operatorname{supp}(\alpha)| \ge (1-\hat{\rho})n.
\]
\end{proposition}

\begin{proof}
Let $u = 1/((1-\hat{\rho})n)$. Since $\hat{\rho}<1$, we have $u \ge 1/n$, so the uniform weight vector is feasible and the set is non-empty. For the support bound, if only $k$ coordinates of $\alpha$ are positive, then
\[
1=\sum_{i=1}^n \alpha_i \le k u = \frac{k}{(1-\hat{\rho})n},
\]
which implies $k \ge (1-\hat{\rho})n$. Therefore CARAT can suppress suspicious clients, but cannot collapse the aggregate onto an arbitrarily small subset of participants.
\end{proof}

\begin{proposition}[Norm bound for the final aggregate]
\label{prop:aggregate-bound}
Suppose each clipped update satisfies $\|\bar{\Delta}_i\|_2 \le \tau_r$ and $\alpha \in \mathcal{W}(\hat{\rho})$. Then the final aggregate
\[
\Delta_{\mathrm{agg}}=\sum_{i=1}^n \alpha_i \bar{\Delta}_i
\]
satisfies
\[
\|\Delta_{\mathrm{agg}}\|_2 \le \tau_r.
\]
\end{proposition}

\begin{proof}
Because $\alpha_i \ge 0$ and $\sum_i \alpha_i=1$, the aggregate is a convex combination of clipped updates. By the triangle inequality,
\[
\|\Delta_{\mathrm{agg}}\|_2
\le
\sum_{i=1}^n \alpha_i \|\bar{\Delta}_i\|_2
\le
\sum_{i=1}^n \alpha_i \tau_r
=
\tau_r.
\]
Thus the final server step inherits the clipping radius as a hard norm upper bound.
\end{proof}

\begin{proposition}[Influence bound from malicious weight mass]
\label{prop:malicious-mass-bound}
Let $\mathcal{M}$ be the malicious clients selected in a round and let
\[
q_{\mathcal{M}}=\sum_{i\in\mathcal{M}}\alpha_i
\]
be their total CARAT weight. Suppose every clipped update satisfies $\|\bar{\Delta}_i\|_2\le\tau_r$. Then the direct malicious contribution obeys
\[
\left\|\sum_{i\in\mathcal{M}}\alpha_i\bar{\Delta}_i\right\|_2
\le q_{\mathcal{M}}\tau_r.
\]
If $q_{\mathcal{M}}<1$ and
\[
\Delta_{\mathcal{B}}=
\sum_{i\notin\mathcal{M}}\frac{\alpha_i}{1-q_{\mathcal{M}}}\bar{\Delta}_i
\]
denotes the benign-only aggregate obtained by renormalizing the benign weights, then
\[
\|\Delta_{\mathrm{agg}}-\Delta_{\mathcal{B}}\|_2
\le 2q_{\mathcal{M}}\tau_r.
\]
\end{proposition}

\begin{proof}
The first inequality follows from the triangle inequality:
\[
\left\|\sum_{i\in\mathcal{M}}\alpha_i\bar{\Delta}_i\right\|_2
\le
\sum_{i\in\mathcal{M}}\alpha_i\|\bar{\Delta}_i\|_2
\le q_{\mathcal{M}}\tau_r.
\]
For the second inequality, write
\[
\Delta_{\mathrm{agg}}=
(1-q_{\mathcal{M}})\Delta_{\mathcal{B}}
+\sum_{i\in\mathcal{M}}\alpha_i\bar{\Delta}_i.
\]
Therefore
\[
\|\Delta_{\mathrm{agg}}-\Delta_{\mathcal{B}}\|_2
\le
q_{\mathcal{M}}\|\Delta_{\mathcal{B}}\|_2
+\left\|\sum_{i\in\mathcal{M}}\alpha_i\bar{\Delta}_i\right\|_2.
\]
Since $\Delta_{\mathcal{B}}$ is a convex combination of clipped benign updates, $\|\Delta_{\mathcal{B}}\|_2\le\tau_r$ by Proposition~\ref{prop:aggregate-bound}. Combining this with the first inequality gives the stated bound.
\end{proof}

\begin{proposition}[Directional invariance of certificate scoring]
\label{prop:directional-invariance}
Suppose two nonzero clipped updates satisfy $\bar{\Delta}_i = a \bar{\Delta}_j$ for some scalar $a>0$. Then their normalized versions in Eq.~\eqref{eq:normalize} coincide, i.e.,
\[
\tilde{\Delta}_i=\tilde{\Delta}_j.
\]
Consequently, their hidden-task certificates are identical:
\[
c_{i,t}=c_{j,t}\qquad \text{for all } t \in \{1,\dots,T\}.
\]
\end{proposition}

\begin{proof}
Eq.~\eqref{eq:normalize} rescales every nonzero clipped update to the common radius $q_r$. Positive collinearity is therefore removed by normalization:
\[
q_r\frac{\bar{\Delta}_i}{\|\bar{\Delta}_i\|_2}
=
q_r\frac{a\bar{\Delta}_j}{\|a\bar{\Delta}_j\|_2}
=
q_r\frac{\bar{\Delta}_j}{\|\bar{\Delta}_j\|_2}.
\]
Substituting the resulting equality into Eq.~\eqref{eq:cert} gives identical task scores. Thus, after clipping, the certificate stage depends on direction rather than residual norm magnitude.
\end{proof}

\begin{proposition}[Task-wise standardization preserves within-task ordering]
\label{prop:task-standardization}
Fix a task $t$ and define a standardized certificate score
\[
\hat{c}_{i,t}=\frac{c_{i,t}-m_t}{s_t+\varepsilon},
\]
where $m_t \in \R$ and $s_t>0$ depend only on task $t$ and not on client $i$. Then for any clients $i,j$,
\[
c_{i,t} \ge c_{j,t}
\quad \Longleftrightarrow \quad
\hat{c}_{i,t} \ge \hat{c}_{j,t}.
\]
Moreover,
\[
\hat{c}_{i,t}-\hat{c}_{j,t}
=
\frac{c_{i,t}-c_{j,t}}{s_t+\varepsilon}.
\]
\end{proposition}

\begin{proof}
For a fixed task $t$, the map $x \mapsto (x-m_t)/(s_t+\varepsilon)$ is affine with strictly positive slope because $s_t+\varepsilon>0$. Hence it preserves ordering. The difference identity follows by direct subtraction. Therefore column-wise robust standardization calibrates cross-task scale without changing the within-task ranking induced by the raw certificate scores.
\end{proof}

\begin{proposition}[Softmin bounds and gradient form]
\label{prop:softmin}
Let $s \in \R^T$ and define
\[
\softmin_{\beta}(s)
= -\frac{1}{\beta}\log\left(\frac{1}{T}\sum_{t=1}^T e^{-\beta s_t}\right).
\]
Then
\[
\min_t s_t \le \softmin_{\beta}(s) \le \min_t s_t + \frac{\log T}{\beta}.
\]
Moreover,
\[
\lim_{\beta \to \infty}\softmin_{\beta}(s)=\min_t s_t,
\qquad
\lim_{\beta \to 0}\softmin_{\beta}(s)=\frac{1}{T}\sum_{t=1}^T s_t.
\]
If $z(\alpha)=C^\top\alpha$ and
\[
\pi_t(\alpha)=\frac{\exp(-\beta z_t(\alpha))}{\sum_{\ell=1}^T \exp(-\beta z_\ell(\alpha))},
\]
then
\[
\nabla_\alpha \softmin_{\beta}(C^\top\alpha)=C\pi(\alpha).
\]
\end{proposition}

\begin{proof}
Let $m=\min_t s_t$. Since
\[
e^{-\beta m}/T \le \frac{1}{T}\sum_{t=1}^T e^{-\beta s_t} \le e^{-\beta m},
\]
applying $-(1/\beta)\log(\cdot)$ yields
\[
m \le \softmin_{\beta}(s) \le m+\frac{\log T}{\beta}.
\]
The limit $\beta \to \infty$ follows from the vanishing gap bound. For $\beta \to 0$, a first-order expansion of the exponential average gives
\[
\frac{1}{T}\sum_{t=1}^T e^{-\beta s_t}
=
1-\beta \frac{1}{T}\sum_{t=1}^T s_t + o(\beta),
\]
so the softmin converges to the arithmetic mean. Differentiating the softmin expression with respect to $z_t$ yields the Gibbs weights $\pi_t(\alpha)$, and the chain rule gives $C\pi(\alpha)$. Hence the ascent direction is a convex combination of certificate columns, with more weight placed on tasks whose current aggregate utility is smaller.
\end{proof}

\begin{proposition}[Temporal prior and projected-iterate invariants]
\label{prop:temporal-prior}
Define
\[
f(\alpha)=\softmin_{\beta}(C^\top\alpha)-\lambda_r \gamma^\top\alpha-\lambda_p\|\alpha-\alpha^{\mathrm{prev}}\|_2^2.
\]
If $\alpha,\alpha' \in \mathcal{W}(\hat{\rho})$ satisfy $C^\top\alpha=C^\top\alpha'$ and $\gamma^\top\alpha=\gamma^\top\alpha'$, then
\[
f(\alpha)-f(\alpha')
=
-\lambda_p\!\left(\|\alpha-\alpha^{\mathrm{prev}}\|_2^2-\|\alpha'-\alpha^{\mathrm{prev}}\|_2^2\right).
\]
In addition, under projected gradient ascent
\[
\alpha^{(k+1)}=\Pi_{\mathcal{W}(\hat{\rho})}\!\left(\alpha^{(k)}+\eta \nabla f(\alpha^{(k)})\right),
\]
every iterate remains in $\mathcal{W}(\hat{\rho})$.
\end{proposition}

\begin{proof}
The first statement follows by direct substitution: if the certificate and rank terms match, the only remaining difference in objective value is the quadratic prior term. CARAT therefore prefers, among utility-equivalent feasible weights, the solution closer to $\alpha^{\mathrm{prev}}$. The second statement is immediate from the projection operator. Hence the optimizer preserves three round invariants throughout: $\alpha_i \ge 0$, $\sum_i \alpha_i =1$, and $\alpha_i \le 1/((1-\hat{\rho})n)$.
\end{proof}

\begin{proposition}[Pairwise weight-transfer condition]
\label{prop:pairwise-transfer}
Let
\[
f(\alpha)=\softmin_{\beta}(C^\top\alpha)-\lambda_r \gamma^\top\alpha-\lambda_p\|\alpha-\alpha^{\mathrm{prev}}\|_2^2
\]
be the CARAT objective, and let $e_a,e_b$ be coordinate vectors for two clients $a$ and $b$. Suppose $\alpha+\epsilon(e_b-e_a)$ remains feasible for sufficiently small $\epsilon>0$. Define
\[
\pi_t(\alpha)=\frac{\exp(-\beta (C^\top\alpha)_t)}
{\sum_{\ell=1}^T \exp(-\beta (C^\top\alpha)_\ell)}.
\]
Then moving infinitesimal weight from client $a$ to client $b$ increases the objective to first order whenever
\[
(C_b-C_a)^\top \pi(\alpha)
-\lambda_r(\gamma_b-\gamma_a)
-2\lambda_p\!\left[
(\alpha_b-\alpha_b^{\mathrm{prev}})
-(\alpha_a-\alpha_a^{\mathrm{prev}})
\right] > 0,
\]
where $C_i$ denotes row $i$ of the standardized certificate matrix.
\end{proposition}

\begin{proof}
Let $d=e_b-e_a$. A first-order expansion gives
\[
f(\alpha+\epsilon d)-f(\alpha)
=\epsilon d^\top\nabla f(\alpha)+o(\epsilon).
\]
Using Proposition~\ref{prop:softmin}, the gradient is
\[
\nabla f(\alpha)=C\pi(\alpha)-\lambda_r \gamma-2\lambda_p(\alpha-\alpha^{\mathrm{prev}}).
\]
Substituting $d=e_b-e_a$ yields exactly the displayed expression. Hence a client gains weight when its hidden-task utility advantage, measured through the current softmin focus $\pi(\alpha)$, outweighs its rank penalty and temporal-prior cost.
\end{proof}

\begin{proposition}[Per-round complexity]
\label{prop:complexity}
Let $n$ be the number of clients, $T$ the number of hidden tasks, $m$ the number of sampled coordinates for rank scoring, $S$ the number of rank subsamples, and $K$ the number of projected-gradient steps. Ignoring constant factors from batch size and model width, the extra defense-side work of CARAT scales as
\[
\text{probe-model forward passes} + O(Snm + KnT).
\]
\end{proposition}

\begin{proof}
CARAT performs one baseline probe evaluation plus $n$ candidate probe evaluations over the $T$ hidden tasks, then computes $S$ rank summaries over $m$ coordinates, and finally executes $K$ projected optimization steps over $n$ weights and $T$ task scores. In the prototype implementation, the probe-model forward passes dominate wall-clock time, while the rank and optimization stages are lower-order terms controlled by $m$, $S$, and $K$.
\end{proof}

\paragraph{Remark on the reference assumption.}
The object $\mathcal{D}^{\mathrm{ref}}$ denotes the labeled server reference set, while $\mathcal{P}$ is a subsampled working pool derived from it for hidden-task construction. This distinction is part of CARAT's threat model and materially strengthens the setting relative to fully data-free robust aggregation.

\subsection{Reproducibility, Ethics, and Availability}
\label{sec:appendix-reproducibility}

\paragraph{Reproducibility.}
All main reported results use the public CIFAR-10 and CIFAR-100 benchmarks and the fixed training protocol described in Section~\ref{sec:experiments}. The aggregation method is specified in Algorithm~\ref{alg:carat-round}, with implementation defaults listed in Appendix~\ref{sec:appendix-experiments}. The main CARAT settings and component variants each use five seeds, and the result archives retain metrics, completion markers, job metadata, and configuration files so that every reported result can be traced to seed-level runs.

\paragraph{Compute resources.}
The main CARAT sweep and reported public-baseline logs were generated through the repository launch scripts on a Compute Canada SLURM cluster; the component study was run with the same experiment interface on a laboratory GPU server. Each run writes a metadata file recording its configuration, command line, seed, Python environment, and CUDA requirement; the main sweep also records the git commit.

\paragraph{Ethics.}
This work studies defenses against model-poisoning attacks in federated learning. The attack implementations are used to evaluate robustness and do not require access to private user data. The main practical risk is dual use: stronger attack simulations may help adversaries tune poisoning strategies. We mitigate this by framing attacks as benchmark stressors and by releasing defensive evaluation code together with the corresponding protocols.

\paragraph{Data availability.}
CIFAR-10 and CIFAR-100 are public datasets. In the reported implementation, CARAT samples a $512$-example server pool from the corresponding public training set used to construct client partitions; those examples are not removed from client local data. This overlap and the unequal CARAT/FLTrust reference budgets are limitations of the current evaluation. A stricter reproduction should additionally test deterministic, disjoint server splits while reporting the resulting reduction in client training data.
